\section{Architecture-First Design Approach}

This project adopts an architecture-first design approach, prioritizing the definition of a
secure and governable system architecture before extensive implementation \cite{sommervilleESP}. This
methodology is particularly suited to security-critical systems, where early architectural
decisions directly influence access control effectiveness, auditability, and threat
mitigation.

Core architectural elements including network segmentation, identity and access
management, privileged access pathways, and monitoring are defined and evaluated at
a high level prior to development. This ensures that design decisions are intentional,
defensible, and aligned with established security principles rather than emerging
organically through feature-driven implementation.

By treating architecture as the primary artifact of Semester One, the project emphasis's
that secure systems are achieved through correct foundational design rather than
retrofitted controls.


\section{Iterative Validation and Prototyping}

While Semester One is design-focused, key architectural assumptions are validated
through limited and targeted prototyping. This allows feasibility and constraints to be
identified without committing to full-scale implementation.

Validation activities include testing logical network segmentation, verifying secure
remote and privileged access pathways, confirming identity-based authentication flows,
and validating centralized log collection. These prototypes are treated as disposable
validation artifacts rather than production components, with outcomes feeding directly
back into architectural refinement. All documentation and technical artefacts are maintained under version control with regular incremental commits.
